\chapter{Taylor's Theorem and Approximation}

In the previous chapter, we built Taylor series $\sum_{k=0}^\infty \frac{f^{(k)}(a)}{k!}(x-a)^k$
from a local-to-global idea: if you know everything about a smooth function --- meaning
a function whose derivatives of all orders exist without kinks or sudden breaks --- at a
\emph{single anchor point} $a$, you can often reconstruct the function far away. Starting from
the tangent line $f(a) + f'(a)(x-a)$, which matches height and slope, adding a second derivative
produces an \emph{osculating parabola} --- from the Latin \emph{osculari}, ``to kiss'' --- that
bends with the exact same curvature (how sharply the curve turns away from its tangent line). A
third derivative buys a cubic matching the rate at which curvature changes, and each further
derivative lets the \emph{Taylor polynomial} $P_n(x)$ hug the curve across a wider window before
peeling away.

In practice, any concrete calculation stops after finitely many terms, and even when we let the
degree grow toward the full \emph{Taylor series}, we saw at the end of Chapter~30 that matching
all derivatives at $a$ does not automatically force the series to equal $f(x)$ away from $a$.
That leaves two central questions:
\begin{enumerate}
  \item \emph{How large is the truncation error}, the error made by cutting the series off (truncating it) after finitely many terms? When approximating $f(x)$ by a Taylor polynomial $P_n(x)$ of degree $n$, how big can the remainder $R_n(x) = f(x) - P_n(x)$ be?
  \item \emph{When does the full series equal $f(x)$?} Under what conditions does $R_n(x) \to 0$ as $n \to \infty$, proving that the Taylor series converges to $f(x)$ rather than drifting away?
\end{enumerate}

Both questions are answered by an explicit formula for the gap $R_n(x)$ at any finite step,
called \emph{Taylor's theorem with remainder}. With it we can bound the error of a Taylor
polynomial with certainty, and see why the series reproduces functions like $e^x$ and
$\cos x$ everywhere but stops working for others.

\section{From a line to a polynomial: matching derivatives}

Start with the linear approximation and recall why it works near $x=a$. Near a point $a$,
the line
\[
  L(x) = f(a) + f'(a)(x-a)
\]
agrees with $f$ in two respects: $L(a) = f(a)$ (same height) and $L'(a) = f'(a)$
(same slope). These two matched conditions are why the line stays close to
the curve near $a$, and it is the only line that gets both right.

\begin{intuitionbox}
Approximation by matching is a recurring technique in mathematics: to model a
complicated function near a point, construct a simple polynomial that agrees with it in as many
\emph{derivatives} as possible. A line matches two pieces of data (value and
slope). A parabola can match three (value, slope, and curvature). A degree-$n$
polynomial has $n+1$ free coefficients, so it can match $n+1$ pieces of data: the
value and the first $n$ derivatives. The more derivatives you match at $a$, the
closer the polynomial approximates the function as you move away from $a$.
\end{intuitionbox}

To construct higher-order approximations, we specify a polynomial $P_n$ of degree $n$ satisfying
\[
  P_n(a) = f(a), \quad P_n'(a) = f'(a), \quad P_n''(a) = f''(a), \quad \dots, \quad
  P_n^{(n)}(a) = f^{(n)}(a).
\]
That is, $P_n$ matches the function and its first $n$ derivatives at the single
point $a$. Which coefficients make this happen?

In Chapter~30 we extracted these coefficients under the optimistic assumption that $f$
already equaled an infinite power series. The same algebra works for a \emph{finite}
polynomial of degree $n$, with no convergence assumption at all. Write $P_n$ in powers of
$(x-a)$, because we only care about behavior near $a$:
\[
  P_n(x) = c_0 + c_1(x-a) + c_2(x-a)^2 + c_3(x-a)^3 + \cdots + c_n(x-a)^n.
\]
Now extract the coefficients one at a time by differentiating and then plugging in
$x=a$, which kills every term that still has a factor of $(x-a)$ in it.

Plug in $x=a$ directly: every term but the first dies, leaving $P_n(a) = c_0$. So
$c_0 = f(a)$.

Differentiate once: $P_n'(x) = c_1 + 2c_2(x-a) + 3c_3(x-a)^2 + \cdots$. Plug in
$x=a$: everything but the first term dies, $P_n'(a) = c_1$. So $c_1 = f'(a)$.

Differentiate again: $P_n''(x) = 2c_2 + 3\cdot 2\, c_3 (x-a) + \cdots$. Plug in
$x=a$: $P_n''(a) = 2c_2$. So $c_2 = f''(a)/2$.

Once more: $P_n'''(x) = 3\cdot 2\, c_3 + 4\cdot 3\cdot 2\, c_4 (x-a) + \cdots$, so
$P_n'''(a) = 6c_3$, giving $c_3 = f'''(a)/6$.

The numbers $1, 2, 6, 24, \dots$ appearing in the denominators are the factorials,
because differentiating $(x-a)^k$ exactly $k$ times brings
down $k \cdot (k-1) \cdots 2 \cdot 1 = k!$ and turns the term into a constant, while
all lower terms have already vanished and all higher terms still carry a factor of
$(x-a)$. So in general
\[
  P_n^{(k)}(a) = k!\, c_k \quad\Longrightarrow\quad c_k = \frac{f^{(k)}(a)}{k!}.
\]

\begin{focusbox}
\textbf{The Taylor polynomial.} The unique degree-$n$ polynomial matching $f$ and its
first $n$ derivatives at $a$ is
\[
  P_n(x) = \sum_{k=0}^{n} \frac{f^{(k)}(a)}{k!}\,(x-a)^k
  = f(a) + f'(a)(x-a) + \frac{f''(a)}{2!}(x-a)^2 + \cdots + \frac{f^{(n)}(a)}{n!}(x-a)^n.
\]
When $a=0$ it is called the \emph{Maclaurin} polynomial, as in Chapter~30.
\end{focusbox}

\begin{whybox}
Why factorials? Each derivative peels one factor off the front of $(x-a)^k$:
$(x-a)^k \to k(x-a)^{k-1} \to k(k-1)(x-a)^{k-2} \to \cdots$. After exactly $k$ steps
the variable is gone and the constant $k!$ is left. Dividing the coefficient by $k!$
in advance cancels it, so when you differentiate $k$ times and set $x=a$, the $k$-th
term contributes exactly $f^{(k)}(a)$ and no other term contributes anything. Each
coefficient then controls one derivative, independently of the others.
\end{whybox}

\begin{recoverybox}
\textbf{Pause and think.}
Before moving on, test the formula on a polynomial: what is the degree-$3$ Taylor polynomial
of $f(x) = 2x^3 - 4x + 7$ centered at $a = 0$? What is its degree-$4$ Taylor polynomial?
If you expand a degree-$n$ polynomial to degree $n$ (or higher) about any center, what must
the remainder be?

\emph{Reflection}: Differentiating a degree-$n$ polynomial $n+1$ times yields zero, so all
derivatives beyond the $n$-th vanish identically. The Taylor polynomial $P_n$ of a degree-$n$
polynomial is the polynomial itself, re-expressed in powers of $(x-a)$. There is no remainder at all.
\end{recoverybox}

Figure~\ref{ch31:fig-matching-layers} shows the first three steps for $f(x)=e^x$ at
$a=0$, where every derivative equals $e^0 = 1$, so $c_k = 1/k!$.

\begin{figure}[H]
\centering
\begin{tikzpicture}
% Three side-by-side panels showing degree 0, 1, 2 approximations to e^x at a=0
\foreach \ly/\shift/\ttl/\poly/\lab/\desc in {%
  0.85/0/{Degree 0: height}/{1}/{$P_0$}/{$P_0(x)=1$\\matches $f(0)$},
  1.35/4.75/{Degree 1: slope}/{1+x}/{$P_1$}/{$P_1(x)=1+x$\\matches $f(0)$, $f'(0)$},
  1.45/9.5/{Degree 2: bending}/{1+x+0.5*x^2}/{$P_2$}/{$P_2(x)=1+x+\tfrac12x^2$\\matches $f(0)$, $f'(0)$, $f''(0)$}} {
\begin{scope}[shift={(\shift,0)}]
  \begin{axis}[
    width=4.9cm, height=4.9cm,
    axis lines=middle,
    xlabel={$x$}, ylabel={$y$},
    xmin=-1.5, xmax=1.5, ymin=-0.2, ymax=3.4,
    samples=100, domain=-1.4:1.4,
    title style={font=\small\bfseries},
    title={\ttl},
    xtick={-1,1}, ytick={1,2,3},
    tick label style={font=\footnotesize},
    clip=true,
  ]
  \addplot[curveblue, very thick] {exp(x)};
  \addplot[tangentred, thick, dashed] {\poly};
  \addplot[only marks, mark=*, mark size=1.8pt, black] coordinates {(0,1)};
  \node[curveblue, font=\footnotesize, anchor=west] at (axis cs:0.75,3.0) {$e^x$};
  \node[tangentred, font=\footnotesize, anchor=north west] at (axis cs:0.5,\ly) {\lab};
  \end{axis}
  \node[below, font=\footnotesize, align=center] at (1.75,-0.25) {\desc};
\end{scope}
}
\end{tikzpicture}
\caption{Building Taylor polynomials of $f(x)=e^x$ at $a=0$ one degree at a time
(blue: $e^x$; red dashed: the polynomial). Degree~0 matches the height, degree~1 adds the
slope (the tangent line), and degree~2 adds the second derivative, so the parabola bends
the same way as the curve. Each new term $x^k/k!$ fixes the $k$-th derivative at $0$
without changing the ones already matched.}
\label{ch31:fig-matching-layers}
\end{figure}
\subsection{A first picture: parabola beats line}
Take $f(x) = \cos x$ at $a = 0$. We know $f(0)=1$, $f'(0) = -\sin 0 = 0$,
$f''(0) = -\cos 0 = -1$. So
\[
  P_1(x) = 1, \qquad P_2(x) = 1 - \tfrac{1}{2}x^2.
\]
The linear approximation of cosine at the top of its hump is the flat line
$y=1$ --- correct in value and slope but blind to the fact that cosine curves
\emph{downward}. The quadratic $1 - \tfrac{1}{2}x^2$ accounts for that downward bend,
and Figure~\ref{ch31:fig-cos-parabola} shows that it follows the cosine much more closely.
\begin{figure}[H]
\centering
\begin{tikzpicture}
\begin{axis}[
  width=10.5cm, height=6.5cm,
  axis lines=middle,
  xlabel={$x$}, ylabel={$y$},
  xmin=-3.3, xmax=3.3, ymin=-1.6, ymax=1.5,
  ytick={-1,-0.5,0,0.5},
  samples=160, domain=-3.1:3.1,
  legend pos=outer north east,
  legend style={font=\small, draw=gray!40, rounded corners=2pt},
  legend cell align=left,
  clip=true,
]
\addplot[curveblue, very thick] {cos(deg(x))};
\addlegendentry{$\cos x$}
\addplot[gray!70!black, thick, dashed, domain=-3.1:3.1] {1};
\addlegendentry{$P_1(x) = 1$}
\addplot[tangentred, thick] {1 - 0.5*x^2};
\addlegendentry{$P_2(x) = 1 - \frac{1}{2}x^2$}
\addplot[only marks, mark=*, mark size=1.8pt, black] coordinates {(0,1)};
\node[above left, font=\footnotesize] at (axis cs:-0.08,1.02) {$(0,1)$};
\end{axis}
\end{tikzpicture}
\caption{At $a=0$, the linear approximation to $\cos x$ is the horizontal tangent line $y=1$,
which matches height and slope but ignores curvature. The quadratic $P_2(x)=1-\tfrac12 x^2$
incorporates the downward bending $f''(0)=-1$ and stays close to the cosine over a much wider interval.}
\label{ch31:fig-cos-parabola}
\end{figure}

Each extra term improves the fit. The degree-$4$ Taylor polynomial of cosine,
$1 - \tfrac12 x^2 + \tfrac{1}{24}x^4$, stays within about $0.08$ of $\cos x$ all the
way out to $x=2$, while $P_2$ is already off by about $0.58$ there. The general
pattern: \emph{higher degree gives a wider interval on which the polynomial can be
trusted.} The rest of the chapter makes ``within'' and ``trusted'' precise.

\section{The remainder: naming the gap exactly}

A picture of a polynomial staying close to a curve suggests a good approximation,
but it proves nothing. The exact
object is the \emph{remainder}, the difference between the function and its
polynomial:
\[
  R_n(x) = f(x) - P_n(x).
\]
This is an exact identity rather than an approximation: $f(x) = P_n(x) + R_n(x)$
for every $x$, with $R_n(x)$ collecting the entire approximation error. Analyzing the approximation now reduces to bounding
$R_n$: if we can show $R_n(x)$ is small, then $P_n(x)$ is a good approximation, with
a guarantee, not a hope.

\begin{happeningbox}
When you write $\cos x \approx 1 - \tfrac12 x^2$, the ``$\approx$'' is hiding a
precise statement. Behind it is an equation with a correction term:
\[
  \cos x = 1 - \tfrac12 x^2 + R_2(x),
\]
and Taylor's theorem is going to hand us a \emph{formula} for $R_2(x)$ --- not its
exact value (that would be asking too much), but an expression that lets us
\emph{bound} it. A bound on this remainder is what separates an engineer who can
justify replacing $\cos x$ by $1 - \tfrac12 x^2$ for small angles from one who only
hopes it works.
\end{happeningbox}

What should we expect of $R_n$? Since $P_n$
matches $f$ and its first $n$ derivatives at $a$, the remainder and \emph{its} first
$n$ derivatives must all vanish at $a$:
\[
  R_n(a) = R_n'(a) = R_n''(a) = \cdots = R_n^{(n)}(a) = 0.
\]
A function whose value and first $n$ derivatives are all zero at a point is very
flat there. So we expect $R_n(x)$ to behave like a constant times $(x-a)^{n+1}$ near
$a$, which shrinks fast as $x \to a$. This is why more terms help near $a$. The next
theorem makes the expectation precise.

\section{Taylor's theorem with the Lagrange remainder}

We state the theorem first and prove it afterwards.

\begin{focusbox}
\textbf{Taylor's Theorem (Lagrange form of the remainder).} Suppose $f$ is $n+1$
times differentiable on an open interval containing $a$ and $x$. Then there is a
point $c$ strictly between $a$ and $x$ such that
\[
  f(x) = \underbrace{\sum_{k=0}^{n} \frac{f^{(k)}(a)}{k!}(x-a)^k}_{P_n(x)}
       \; + \; \underbrace{\frac{f^{(n+1)}(c)}{(n+1)!}\,(x-a)^{n+1}}_{R_n(x)}.
\]
The remainder $R_n(x)$ looks exactly like the \emph{next} Taylor term --- the
degree-$(n+1)$ term you would have written next --- except that the
$(n+1)$-st derivative is evaluated not at $a$ but at some unknown point $c$ in between.
\end{focusbox}

This way of writing the remainder is named after Joseph-Louis Lagrange, who found it
(see the history box below).

This generalizes the Mean Value Theorem (MVT) of Chapter~15. Set $n=0$. The
theorem says $f(x) = f(a) + f'(c)(x-a)$ for some $c$ between $a$ and $x$ --- which,
rearranged, is $\frac{f(x)-f(a)}{x-a} = f'(c)$, precisely the MVT. So Taylor's
theorem is the MVT with more derivatives. The point $c$ where the MVT's average
slope is achieved reappears here as the point where the
$(n+1)$-st derivative takes the value that makes the equation exactly balance.

\begin{intuitionbox}
A natural first guess is that the truncation error after $n$ terms should look like
the first omitted term, $\frac{f^{(n+1)}(a)}{(n+1)!}(x-a)^{n+1}$.
Taylor's theorem says the error has exactly that form, except that the derivative is
evaluated at a shifted point $c$ instead of at $a$. Since you usually need only a
\emph{bound} on $f^{(n+1)}$ over the interval, not $c$ itself, the shift costs little
in practice. You give up knowing the exact error (impossible) and get the exact
\emph{form} of the error (immediately useful).
\end{intuitionbox}

\begin{historybox}
Between the fourteenth and nineteenth centuries, calculus changed from geometric
intuition into numerical calculation with guaranteed precision.
The earliest origins trace to Madhava of Sangamagrama and the Kerala school in fourteenth-century
India, who discovered the infinite series expansions for $\sin x$, $\cos x$, and $\arctan x$
to construct accurate astronomical tables. Centuries later in Europe, James Gregory (in 1671)
and Isaac Newton independently rediscovered several of these expansions. In 1715, the English
mathematician Brook Taylor published the general recipe matching derivatives to polynomial
coefficients in his \emph{Methodus Incrementorum Directa et Inversa}.

Yet a piece was missing: \emph{none of them had the remainder}.
They had no formula for the error committed by truncating a series after a finite number of terms.
The early calculus of Newton, Leibniz, and Euler treated infinite series as formal algebraic
identities. Mathematicians calculated with them boldly and often correctly by instinct, but they
were working without safety nets. When Leonhard Euler manipulated divergent series --- infinite
sums whose running totals oscillate or blow up toward infinity rather than settling to a fixed
number --- to deduce surprising identities, other mathematicians grew uneasy,
unable to draw a clear line between sound approximation and mathematical illusion.

Controversy simmered on two fronts. Taylor and Johann Bernoulli disputed who had first
discovered the polynomial expansion, while Colin Maclaurin's 1742 treatise
popularized the expansion centered at zero so widely that history attached Maclaurin's name
to Taylor's special case. More fundamentally, mathematicians debated whether infinite sums could
ever be trusted without an error term.

The answer came in 1797, when Joseph-Louis Lagrange published \emph{Th{\'e}orie des
fonctions analytiques}. There he isolated the exact, finite remainder
term that bears his name. In the 1820s, Augustin-Louis Cauchy and later Karl Weierstrass placed
Lagrange's remainder on rigorous footing by linking it directly to the Mean Value Theorem.

With a computable remainder bound, calculus became an engineering tool:
navigators could rely on printed tables of logarithms,
astronomers could compute orbits with certified precision, and today computer chips
evaluate trigonometric functions knowing the exact number of terms needed to make every displayed
digit correct. The one-variable Taylor remainder is fully settled, but finding sharp,
computable error bounds for approximations involving many variables, such as those used in
machine learning, is still an active area of research.
\end{historybox}

\subsection{Why it is true: the idea before the proof}

There are several routes to this theorem (an exercise develops one through repeated
integration by parts). Here we take a direct route through the Cauchy Mean Value
Theorem of Chapter~15. It says: if $g$ and $h$ are continuous on the closed interval
between $a$ and $b$, differentiable inside it, and $h'$ is never zero inside, then
\[
  \frac{g(a) - g(b)}{h(a) - h(b)} = \frac{g'(c)}{h'(c)}
\]
for some $c$ strictly between $a$ and $b$.

We want to compare $f(x)$ with $P_n(x)$. Fix the target point $x = b$ and the center
$a$. Build a helper function $g(t)$ of a new variable $t$ that measures, at each $t$,
how far $f$ is from \emph{its own Taylor polynomial centered at $t$} (evaluated at
$b$), and arrange it so that $g$ equals $0$ at the end $t=b$. When we differentiate
$g$, the sum telescopes (each term cancels part of its neighbor, as in Chapter~21, so
only one term survives), and one application of the Cauchy Mean Value Theorem then
gives the remainder formula.

\begin{whybox}
Why a helper function? Because the quantity we care about, $f(b) - P_n(b)$, is a
single number with no $t$ in it --- there is nothing to apply a mean-value argument
to. The trick, used throughout analysis (the rigorous theory behind calculus), is to
\emph{introduce a variable}: turn the fixed number into the endpoint value of a
moving quantity, so that theorems about derivatives, such as Rolle's theorem and the
MVT, can be applied. We let the fixed center $a$ slide to a variable center $t$,
watch the discrepancy change as $t$ moves from $a$ to $b$, and let the Mean Value
Theorem find a point where the rate of change reveals the remainder. Choosing the
right helper function is the hard step; checking that it works is then routine.
\end{whybox}

\subsection{The proof}

\begin{theorem}[Taylor with Lagrange remainder]
\label{ch31:thm-taylor}
Let $f$ be a function such that $f, f', \dots, f^{(n)}$ are continuous on the closed
interval between $a$ and $b$, and $f^{(n+1)}$ exists on the open interval between
them. Then there is a point $c$ strictly between $a$ and $b$ with
\[
  f(b) = \sum_{k=0}^{n} \frac{f^{(k)}(a)}{k!}(b-a)^k
       + \frac{f^{(n+1)}(c)}{(n+1)!}(b-a)^{n+1}.
\]
\end{theorem}

\begin{proof}
Treat $a$, $b$, $n$ as fixed. Define a function of a sliding center $t$, living on
the interval between $a$ and $b$:
\[
  g(t) = f(b) - \sum_{k=0}^{n} \frac{f^{(k)}(t)}{k!}\,(b-t)^k.
\]
In words, $g(t)$ is the gap between the true value $f(b)$ and the Taylor
polynomial of $f$ \emph{centered at $t$}, evaluated at $b$. Two endpoint values are
immediate.

At $t = b$: every factor $(b-t)^k$ with $k \ge 1$ vanishes, and the $k=0$ term is
$f(b)$, so
\[
  g(b) = f(b) - f(b) = 0.
\]
At $t = a$: the sum is exactly the Taylor polynomial $P_n(b)$, so
\[
  g(a) = f(b) - P_n(b) = R_n(b),
\]
which is the remainder $R_n(b)$ we wish to evaluate. Thus $g$ evaluates to $R_n(b)$
at $t=a$ and to $0$ at $t=b$.

Now differentiate $g$ with respect to $t$. Each term of the sum is a product $\frac{f^{(k)}(t)}{k!}(b-t)^k$, so we use the
product rule:
\[
  \frac{\dd}{\dd t}\!\left[\frac{f^{(k)}(t)}{k!}(b-t)^k\right]
  = \frac{f^{(k+1)}(t)}{k!}(b-t)^k
    \;-\; \frac{f^{(k)}(t)}{k!}\,k\,(b-t)^{k-1}.
\]
The first piece is the derivative landing on $f^{(k)}$; the second is the derivative
landing on $(b-t)^k$, which brings down $k$ and lowers the power (the minus sign
comes from the chain rule on $b-t$). Now sum over $k = 0$ to $n$: the
second pieces cancel the first pieces of the neighboring term. For $k=0$ the term is
$f(t)$, whose derivative is $A_0 = f'(t)$, with no second piece at all. For each
$k\ge 1$, write the first piece as $A_k = \frac{f^{(k+1)}(t)}{k!}(b-t)^k$ and simplify
$\frac{k}{k!} = \frac{1}{(k-1)!}$ in the second piece, which becomes
$\frac{f^{(k)}(t)}{(k-1)!}(b-t)^{k-1} = A_{k-1}$. So
\begin{align*}
  g'(t) &= -A_0 - \sum_{k=1}^{n}\big(A_k - A_{k-1}\big) \\
        &= -\sum_{k=0}^{n} A_k + \sum_{k=0}^{n-1} A_k \\
        &= -A_n,
\end{align*}
because every neighboring pair cancels, leaving only the last first-type term.
Therefore
\[
  g'(t) = -\frac{f^{(n+1)}(t)}{n!}\,(b-t)^n.
\]

To turn this into the remainder formula, we compare $g$ with a second helper function,
chosen so that the $(b-t)^n$ factor will cancel:
\[
  h(t) = (b-t)^{n+1}.
\]
Both $g$ and $h$ are continuous on the closed interval and differentiable on the open
one, and $h'(t) = -(n+1)(b-t)^n$ never vanishes inside. The Cauchy Mean
Value Theorem then guarantees a point $c$ strictly between $a$ and $b$ with
\[
  \frac{g(a) - g(b)}{h(a) - h(b)} = \frac{g'(c)}{h'(c)}.
\]
We computed $g(b) = 0$, $g(a) = R_n(b)$, $h(b) = 0$, and $h(a) = (b-a)^{n+1}$. The
right side is
\[
  \frac{g'(c)}{h'(c)}
  = \frac{-\frac{f^{(n+1)}(c)}{n!}(b-c)^n}{-(n+1)(b-c)^n}
  = \frac{f^{(n+1)}(c)}{(n+1)!}.
\]
The $(b-c)^n$ factors cancel (they are not zero, because $c \ne b$), and
$n!\,(n+1) = (n+1)!$. Setting the two sides equal,
\[
  \frac{R_n(b)}{(b-a)^{n+1}} = \frac{f^{(n+1)}(c)}{(n+1)!},
\]
which rearranges to
\[
  R_n(b) = \frac{f^{(n+1)}(c)}{(n+1)!}(b-a)^{n+1}.
\]
That is the claim.
\end{proof}

\begin{remark}
The case $n=0$ recovers the Mean Value Theorem itself, and the proof above, stripped
to $n=0$, \emph{is} the standard proof of the MVT via the Cauchy form.
Figure~\ref{ch31:fig-sliding-center} summarizes the proof.
\end{remark}

\begin{figure}[H]
\centering
\begin{tikzpicture}
% Sliding center t from a to b
\draw[->, thick, >=stealth] (-0.5,0) -- (8.5,0) node[right, font=\small] {$t$};
\fill (1,0) circle (2pt) node[below=3pt, font=\small] {$a$};
\fill (7,0) circle (2pt) node[below=3pt, font=\small] {$b$};
\fill[quotientorange] (4.5,0) circle (2pt) node[below left=1pt, font=\small, quotientorange] {$c$};

% Top bracket showing interval
\draw[thick, decoration={brace, amplitude=6pt}, decorate] (1,0.3) -- (7,0.3)
  node[midway, above=8pt, font=\small] {Sliding center $t \in [a,b]$};

% Two callout boxes for endpoints
\node[align=center, draw=sourceblue!60, fill=sourceblue!8, rounded corners=2pt, font=\small, inner sep=4pt]
  at (1, 1.8) {At $t=a$:\\$g(a) = f(b) - P_n(b) = R_n(b)$\\(the remainder we want)};

\draw[->, >=stealth, sourceblue, shorten >=2pt] (1, 1.1) -- (1, 0.4);

\node[align=center, draw=targetgreen!60, fill=targetgreen!8, rounded corners=2pt, font=\small, inner sep=4pt]
  at (7, 1.8) {At $t=b$:\\$g(b) = f(b) - f(b) = 0$\\(zero)};

\draw[->, >=stealth, targetgreen, shorten >=2pt] (7, 1.1) -- (7, 0.4);

\node[align=center, draw=quotientorange!80, fill=quotientorange!10, rounded corners=2pt, font=\small, inner sep=4pt]
  at (4.5, -1.3) {Cauchy Mean Value Theorem:\\$\dfrac{g(a)-g(b)}{h(a)-h(b)} = \dfrac{g'(c)}{h'(c)} \implies R_n(b) = \dfrac{f^{(n+1)}(c)}{(n+1)!}(b-a)^{n+1}$};

\draw[->, >=stealth, quotientorange, shorten >=2pt] (4.5, -0.6) -- (4.5, -0.05);
\end{tikzpicture}
\caption{The sliding-center strategy in the proof of Taylor's theorem. The helper function $g(t)$
measures the discrepancy at $b$ when the expansion is centered at $t$. Moving $t$ from $a$ to $b$
connects the unknown remainder $R_n(b)$ at one end to zero at the other. The Cauchy Mean
Value Theorem, applied to $g$ and $h(t)=(b-t)^{n+1}$, supplies a point $c$ in between, and the
telescoped formula for $g'$ turns this into the Lagrange remainder.}
\label{ch31:fig-sliding-center}
\end{figure}

\section{Using the remainder to bound error}

The point $c$ is unknown, and that is fine, because we almost never want $c$ itself.
We want a guarantee, and a guarantee follows as soon as we can say how big
$f^{(n+1)}$ can be on the interval. This is the form of Taylor's theorem used in
practice.

\begin{focusbox}
\textbf{The error bound.} If $\abs{f^{(n+1)}(t)} \le M$ for all $t$ between $a$ and
$x$, then
\[
  \abs{R_n(x)} \;\le\; \frac{M}{(n+1)!}\,\abs{x-a}^{n+1}.
\]
This follows from the Lagrange form, because $c$ lies between $a$ and $x$, so
$\abs{f^{(n+1)}(c)} \le M$. You do not need to find $c$; you only need a ceiling $M$
on the next derivative.
\end{focusbox}

To approximate $f(x)$ by $P_n(x)$ and \emph{know}
how wrong you might be, you ask one question: how large can the $(n+1)$-st derivative
get on the way from $a$ to $x$? Whatever that ceiling $M$ is, the error cannot exceed
$\frac{M}{(n+1)!}\abs{x-a}^{n+1}$. Two competing factors appear in that bound:
$\abs{x-a}^{n+1}$ increases with distance from $a$, while $(n+1)!$ in the denominator
grows rapidly with $n$. For functions whose derivatives stay bounded (like $\sin$, $\cos$,
$e^x$ on any fixed interval), the factorial grows faster, and the bound goes to
zero as you take more terms.

\begin{recoverybox}
\textbf{Pause and think.}
Look at the error bound $\frac{M}{(n+1)!}\abs{x-a}^{n+1}$.
If you double your distance from the base point $a$ (say from $\abs{x-a}=0.1$ to $0.2$),
by what factor does the worst-case error bound grow for a tangent line ($n=1$)?
What about for a degree-$4$ approximation ($n=4$)?

\emph{Reflection}: For $n=1$, the bound scales as $2^{1+1} = 2^2 = 4$. For $n=4$, it scales
as $2^{4+1} = 2^5 = 32$. Higher-degree polynomials provide much tighter bounds near the center,
but the error bound grows more rapidly with distance.
\end{recoverybox}

\begin{labbox}
Approximate $\cos(0.5)$ using $P_2(x) = 1 - \tfrac12 x^2$ and bound the error
\emph{before} comparing to a calculator.

The estimate: $P_2(0.5) = 1 - \tfrac12(0.25) = 1 - 0.125 = 0.875$.

The bound: here $n=2$, so we need a ceiling on $\abs{f^{(3)}}$. For $f=\cos$,
$f'''(x) = \sin x$, and $\abs{\sin t} \le 1$ everywhere, so $M=1$. Then
\[
  \abs{R_2(0.5)} \le \frac{1}{3!}\,(0.5)^3 = \frac{0.125}{6} \approx 0.0208.
\]
So we may assert, with certainty, that the true $\cos(0.5)$ lies in
$[0.875 - 0.021,\; 0.875 + 0.021] = [0.854,\, 0.896]$.

The check: a calculator gives $\cos(0.5) = 0.87758\ldots$, and indeed
$\abs{0.87758 - 0.875} = 0.0026$, comfortably under our promised ceiling of
$0.0208$. The bound held, and it came out larger than the actual error, as a safe
guarantee should. Figure~\ref{ch31:fig-cos-error} compares the actual error with the
bound for every $x$ from $0$ to $1.2$.
\end{labbox}

\begin{figure}[H]
\centering
\begin{tikzpicture}
\begin{axis}[
  width=11cm, height=6.5cm,
  axis lines=left,
  xlabel={$x$}, ylabel={error},
  xmin=0, xmax=1.25, ymin=0, ymax=0.3,
  xtick={0,0.25,0.5,0.75,1},
  ytick={0,0.1,0.2,0.3},
  yticklabel style={/pgf/number format/fixed, /pgf/number format/precision=1},
  tick label style={font=\footnotesize},
  samples=120, domain=0:1.2,
  legend pos=north west, legend cell align=left,
  legend style={font=\small, draw=gray!40, rounded corners=2pt},
  clip=true,
]
\addplot[tangentred, thick, dashed] {x^3/6};
\addlegendentry{bound $\abs{x}^3/3!$ (with $M=1$)}
\addplot[curveblue, very thick] {abs(cos(deg(x)) - (1 - 0.5*x^2))};
\addlegendentry{actual error $\abs{\cos x - P_2(x)}$}
\addplot[only marks, mark=*, mark size=1.6pt, tangentred] coordinates {(0.5,0.0208)};
\addplot[only marks, mark=*, mark size=1.6pt, curveblue] coordinates {(0.5,0.0026)};
\end{axis}
\end{tikzpicture}
\caption{Error of the approximation $\cos x \approx 1 - \tfrac12 x^2$. The red dashed curve
is the guaranteed ceiling $x^3/6$ from the error bound; the blue curve is the actual error.
The actual error always stays below the ceiling. At $x = 0.5$ (dots) the ceiling is $0.0208$
and the actual error is $0.0026$.}
\label{ch31:fig-cos-error}
\end{figure}

\begin{trapbox}
A bound is an upper limit, not a prediction. Beginners sometimes report ``the error
is $0.0208$'' when the theorem only promises ``the error is \emph{at most}
$0.0208$.'' The actual error here was $0.0026$, almost ten times smaller. The
Lagrange bound is deliberately conservative: it assumes the next derivative sits at
its worst-case value $M$ over the whole interval, which it usually does not. Never
quote the bound as the error; quote it as the worst case.

There is another reason the actual error was so small: for cosine, the third
derivative at the origin is $-\sin(0) = 0$. That means the quadratic polynomial
$P_2(x) = 1 - \tfrac12 x^2$ is \emph{also} the cubic Taylor polynomial $P_3(x)$.
If we apply Taylor's theorem at degree $n=3$, we need a ceiling $M_4$ on
$\abs{f^{(4)}} = \abs{\cos}$, and $M_4 = 1$ works, so the remainder is bounded by
\[
  \frac{M_4}{4!}(0.5)^4 = \frac{1}{24}(0.0625) \approx 0.00260,
\]
almost exactly the true error $0.0026$. Whenever the next coefficient happens to be
zero, your polynomial automatically gains an extra degree of accuracy.
\end{trapbox}

\begin{labbox}
A harder one: how many terms of the Maclaurin series for $e^x$ do
we need to compute $e = e^1$ to within $10^{-6}$?

For $f(x) = e^x$ centered at $a=0$, every derivative is $e^x$ again, and on
$[0,1]$ we have $e^t \le e < 3$, so we may take $M = 3$. With $x=1$,
\[
  \abs{R_n(1)} \le \frac{3}{(n+1)!}\,\cdot 1^{n+1} = \frac{3}{(n+1)!}.
\]
We need $\frac{3}{(n+1)!} \le 10^{-6}$, that is, $(n+1)! \ge 3\times 10^{6}$. Now list the
factorials: $9! = 362880$ (too small), $10! = 3628800 \ge 3\times10^6$.
So $(n+1) = 10$, meaning $n = 9$: the polynomial through the $x^9/9!$ term suffices.
Computing $\sum_{k=0}^{9} \frac{1}{k!}$ gives $2.7182815\ldots$ against the true
$e = 2.7182818\ldots$, an error of about $3\times 10^{-7}$ --- under our budget, as
promised. Ten terms ($k = 0$ to $9$) already give six correct decimal places,
because the factorial in the denominator grows so fast
(Figure~\ref{ch31:fig-e-terms}).
\end{labbox}

\begin{figure}[H]
\centering
\begin{tikzpicture}
\begin{axis}[
  width=11cm, height=6.5cm,
  ymode=log,
  xlabel={degree $n$ of the Taylor polynomial},
  ylabel={error bound $3/(n+1)!$},
  xmin=-0.5, xmax=11.5, ymin=1e-9, ymax=10,
  xtick={0,1,...,11},
  ytick={1e-8,1e-6,1e-4,1e-2,1},
  tick label style={font=\footnotesize},
  label style={font=\small},
  clip=true,
]
\addplot[curveblue, thick, mark=*, mark size=2pt] coordinates {
  (0,3) (1,1.5) (2,0.5) (3,0.125) (4,0.025) (5,0.0041667)
  (6,0.00059524) (7,7.4405e-5) (8,8.2672e-6) (9,8.2672e-7)
  (10,7.5157e-8) (11,6.2631e-9)};
\draw[trapred, thick, dashed] (axis cs:-0.5,1e-6) -- (axis cs:11.5,1e-6);
\addplot[only marks, mark=*, mark size=3pt, trapred] coordinates {(9,8.2672e-7)};
\node[trapred, font=\small, anchor=north west] at (axis cs:-0.3,1e-6) {target $10^{-6}$};
\node[trapred, font=\small, anchor=west] at (axis cs:9.2,2e-6) {$n=9$};
\end{axis}
\end{tikzpicture}
\caption{The error bound $3/(n+1)!$ for approximating $e$ by the degree-$n$ Taylor
polynomial of $e^x$, on a logarithmic scale (each labeled step on the vertical axis is a factor of $100$).
The first point below the red target line $10^{-6}$ is $n = 9$.}
\label{ch31:fig-e-terms}
\end{figure}

\begin{applicationbox}
Software libraries evaluate $e^x$, $\sin x$, and related functions in a similar way.
First they move the input into a small standard interval near the center, using identities
such as $\sin(x + 2\pi) = \sin x$ (this step is called \emph{range reduction}). Then they
compute a polynomial of fixed degree, and a remainder bound of the kind above guarantees
that the result is accurate to machine precision (the roughly $16$ significant digits a
computer stores). In practice the polynomial is often a slightly adjusted version of the
Taylor polynomial, tuned to spread the error evenly over the interval, but the error
analysis follows the same idea.
\end{applicationbox}

\section{When does the series converge to the function?}

Suppose $f$ has \emph{all} derivatives. We can keep building
Taylor polynomials of every degree, and we may ask whether the infinite series
\[
  \sum_{k=0}^{\infty} \frac{f^{(k)}(a)}{k!}(x-a)^k
\]
adds up to $f(x)$. The series equals $f(x)$ at a given point $x$ exactly when the
remainder vanishes in the limit:
\[
  \boxed{\;\lim_{n\to\infty} R_n(x) = 0\;}
\]
This equation is the exact criterion for a Taylor series to represent its function.
The Taylor series converges \emph{to $f$} at $x$ if and only if the leftover $R_n(x)$ shrinks to
zero as you include more terms. Convergence of the series alone is not enough: as
we will see, a Taylor series can converge to the \emph{wrong} function.

\begin{whybox}
Why is $R_n \to 0$ the exact condition, and not something about the series
``converging''? Because $f(x) = P_n(x) + R_n(x)$ is an identity for every $n$.
Rearranged, $R_n(x) = f(x) - P_n(x)$. As $n \to \infty$, $P_n(x)$ is the partial sum
of the Taylor series. So $P_n(x) \to f(x)$ --- the series summing to the function ---
is the same statement as $R_n(x) \to 0$. That is why this chapter concentrates on
the remainder.
\end{whybox}

For $\sin x$, $\cos x$ and $e^x$, the error bound settles the question quickly. Take
$\sin x$ centered at $0$. Every derivative is one of $\pm\sin x, \pm\cos x$, so
$\abs{f^{(n+1)}(t)} \le 1$ for all $t$, giving
\[
  \abs{R_n(x)} \le \frac{1}{(n+1)!}\,\abs{x}^{n+1}.
\]
For any fixed $x$, no matter how large, $\frac{\abs{x}^{n+1}}{(n+1)!} \to 0$ as
$n\to\infty$, because the factorial eventually outgrows any fixed power (Chapter~28).
Therefore the Maclaurin
series for $\sin x$ converges to $\sin x$ for \emph{every} real $x$. The same
argument works for $e^x$ at any fixed $x$: every derivative is $e^t$, which between
$0$ and $x$ is at most $e^x$ when $x\ge 0$ and at most $e^0 = 1$ when $x<0$. Either
way $M = e^{\abs{x}}$ bounds $\abs{f^{(n+1)}(t)}$ independently of $n$, so
$\abs{R_n(x)} \le e^{\abs{x}}\frac{\abs{x}^{n+1}}{(n+1)!} \to 0$. These functions are
reconstructed everywhere from their derivatives at the single point $x=0$.
Figure~\ref{ch31:fig-sin-taylor} shows this happening for $\sin x$.

\begin{figure}[H]
\centering
\begin{tikzpicture}
\begin{axis}[
  width=11cm, height=7cm,
  axis lines=middle,
  xlabel={$x$}, ylabel={$y$},
  xmin=-7, xmax=7, ymin=-2.2, ymax=2.2,
  xtick={-4,-2,2,4},
  samples=200, domain=-6.8:6.8,
  legend pos=outer north east, legend cell align=left,
  legend style={font=\small, draw=gray!40, rounded corners=2pt},
  restrict y to domain=-2.4:2.4,
  clip=true,
]
\addplot[curveblue, very thick] {sin(deg(x))};
\addlegendentry{$\sin x$}
\addplot[gray, thick, dashed] {x - x^3/6};
\addlegendentry{degree $3$}
\addplot[quotientorange, thick] {x - x^3/6 + x^5/120 - x^7/5040};
\addlegendentry{degree $7$}
\addplot[tangentred, thick] {x - x^3/6 + x^5/120 - x^7/5040 + x^9/362880 - x^11/39916800};
\addlegendentry{degree $11$}
\end{axis}
\end{tikzpicture}
\caption{Taylor polynomials of $\sin x$ at $0$ closing in on the curve. Each higher
degree extends the interval over which the polynomial tracks $\sin x$ before it finally
shoots off to infinity. In the limit, the series captures $\sin x$ on the entire real
line.}
\label{ch31:fig-sin-taylor}
\end{figure}

\subsection{Where convergence stops: the geometric warning}

Not every smooth function is captured on the whole line. Recall from the end of
Chapter~30 the bell-shaped curve
\[
  f(x) = \frac{1}{1+x^2}.
\]
This function is perfectly smooth everywhere on $\R$ --- no breaks, no corners, no
vertical asymptotes, infinitely differentiable. The finite geometric sum
$1 + r + r^2 + \cdots + r^n = \frac{1 - r^{n+1}}{1-r}$ with $r = -x^2$ gives
\[
  \frac{1}{1+x^2} = 1 - x^2 + x^4 - \cdots + (-1)^n x^{2n} + \frac{(-1)^{n+1}x^{2n+2}}{1+x^2},
\]
so the exact remainder after degree $2n$ is
\[
  R_{2n}(x) = \frac{(-1)^{n+1}x^{2n+2}}{1+x^2}.
\]
For $\abs{x} < 1$ it shrinks to $0$; for $\abs{x} > 1$ it grows without bound. So the series
\[
  1 - x^2 + x^4 - x^6 + \cdots
\]
converges to $f$ only for $\abs{x} < 1$. Why $1$? There
is nothing visibly wrong with the function at $x = 1$: as
Figure~\ref{ch31:fig-geometric} shows, the graph passes smoothly through
$(1, \tfrac12)$.

\begin{recoverybox}
\textbf{Pause and think.}
The bell-shaped curve $f(x) = \frac{1}{1+x^2}$ is bounded, completely smooth, and shows no
break or corner anywhere on the real line. If you were confined to the real axis with no
knowledge of complex numbers, could you find any clue in the real graph explaining why its
Taylor series at $a=0$ refuses to converge when $\abs{x} > 1$?

\emph{Reflection}: Nothing in the real graph gives it away. The curve stays bounded and
smooth on $\R$, yet for $\abs{x} > 1$ the partial sums oscillate with growing amplitude. The
obstruction is invisible until you look off the real axis into the imaginary direction.
\end{recoverybox}

\begin{figure}[H]
\centering
\begin{tikzpicture}
\begin{axis}[
  width=10cm, height=6.5cm,
  axis lines=middle,
  xlabel={$x$}, ylabel={$y$},
  xmin=-2.2, xmax=2.2, ymin=-0.6, ymax=1.5,
  ytick={0.5, 1},
  xticklabel style={fill=white, inner sep=1pt, font=\footnotesize},
  yticklabel style={fill=white, inner sep=1pt, font=\footnotesize},
  samples=160, domain=-2.1:2.1,
  legend pos=outer north east,
  legend style={font=\small, draw=gray!40, rounded corners=2pt},
  legend cell align=left,
  restrict y to domain=-0.8:1.7,
  clip=true,
]
\addplot[curveblue, very thick] {1/(1+x^2)};
\addlegendentry{$\frac{1}{1+x^2}$}
\addplot[quotientorange, thick, dashed, domain=-1.4:1.4] {1 - x^2 + x^4 - x^6 + x^8};
\addlegendentry{partial sum (deg $8$)}
\addplot[only marks, mark=*, mark size=1.4pt, black] coordinates {(1,0.5) (-1,0.5)};
\draw[trapred, thick, dotted] (axis cs:1,-0.6) -- (axis cs:1,1.5);
\node[trapred, font=\footnotesize, anchor=west] at (axis cs:1.05,1.3) {$x=1$};
\draw[trapred, thick, dotted] (axis cs:-1,-0.6) -- (axis cs:-1,1.5);
\node[trapred, font=\footnotesize, anchor=east] at (axis cs:-1.05,1.3) {$x=-1$};
\end{axis}
\end{tikzpicture}
\caption{The function $\frac{1}{1+x^2}$ is smooth everywhere, yet its Maclaurin series
$1-x^2+x^4-\cdots$ tracks it only on $(-1,1)$ (dotted lines). Inside that interval the
partial sums close in; just outside, they fly apart, even though the function itself is
perfectly well behaved there. The wall at $\abs{x}=1$ is set by a singularity hiding off
the real axis.}
\label{ch31:fig-geometric}
\end{figure}

\begin{intuitionbox}
The explanation comes from complex numbers (Chapter~30), an insight of
nineteenth-century analysis (Cauchy, then Weierstrass and Riemann). The function $\frac{1}{1+x^2}$, viewed over the complex
plane, \emph{blows up} at $x = \pm i$, where $1 + x^2 = 0$. Those poles sit at
distance $1$ from the origin. A power series centered at $0$ converges inside the
largest disk around $0$ that reaches no singularity --- and the nearest singularities,
invisible on the real line, are exactly one unit away in the imaginary direction. The
radius of convergence $1$ was being dictated by points you could not see.
Figure~\ref{ch31:fig-complex-poles} draws this disk.
\end{intuitionbox}

\begin{figure}[H]
\centering
\begin{tikzpicture}
\draw[->, thick, >=stealth] (-3.2,0) -- (3.2,0) node[right, font=\small] {$\mathrm{Re}(z)$};
\draw[->, thick, >=stealth] (0,-2.6) -- (0,2.6) node[above, font=\small] {$\mathrm{Im}(z)$};

\fill[sourceblue!10] (0,0) circle (1.8cm);
\draw[sourceblue, thick, dashed] (0,0) circle (1.8cm);

\draw[targetgreen, line width=2.4pt] (-1.8,0) -- (1.8,0);

\fill[black] (0,0) circle (2pt) node[below left=2pt, font=\footnotesize] {$0$};

\draw[trapred, very thick] (-0.12, 1.8-0.12) -- (0.12, 1.8+0.12);
\draw[trapred, very thick] (-0.12, 1.8+0.12) -- (0.12, 1.8-0.12);
\node[trapred, font=\small, anchor=west] at (0.3, 2.1) {Pole at $z = +i$};

\draw[trapred, very thick] (-0.12, -1.8-0.12) -- (0.12, -1.8+0.12);
\draw[trapred, very thick] (-0.12, -1.8+0.12) -- (0.12, -1.8-0.12);
\node[trapred, font=\small, anchor=west] at (0.3, -2.1) {Pole at $z = -i$};

\fill[targetgreen!60!black] (1.8,0) circle (2pt) node[below right=1pt, font=\footnotesize] {$1$};
\fill[targetgreen!60!black] (-1.8,0) circle (2pt) node[below left=1pt, font=\footnotesize] {$-1$};

\draw[<->, >=stealth, quotientorange, thick] (0,0) -- (1.27, 1.27)
  node[midway, above left=-1pt, font=\footnotesize, quotientorange] {$R = 1$};

\node[sourceblue!80!black, font=\footnotesize, align=center] at (-2.4, 1.9) {Disk of convergence\\$\abs{z} < 1$};
\draw[->, >=stealth, sourceblue!80!black, shorten >=2pt] (-2.0, 1.6) -- (-1.3, 1.1);

\node[targetgreen!60!black, font=\footnotesize, align=center] at (-3.0, -1.6) {Real interval $(-1,1)$\\visible on $\R$};
\draw[->, >=stealth, targetgreen!60!black, shorten >=2pt] (-2.6, -1.2) -- (-0.9, -0.08);
\end{tikzpicture}
\caption{Why the Maclaurin series of $\frac{1}{1+x^2}$ stops at $\abs{x}=1$. Writing $z$ for a
complex input, the function blows up at $z = \pm i$ (red crosses), where the denominator $1+z^2$
is zero. A power series centered at $0$ converges inside the largest circle around $0$ that
contains no singularity; here that circle has radius $R=1$. The interval of convergence $(-1,1)$
(green) is where this disk meets the real axis.}
\label{ch31:fig-complex-poles}
\end{figure}

\begin{connectionbox}
This connects to complex analysis, the study of calculus for functions of a complex
variable. A real power series centered at $a$ converges
on an interval $(a-R, a+R)$ (the green slice in Figure~\ref{ch31:fig-complex-poles});
the number $R$ is the \emph{radius of convergence}, and
it equals the distance from $a$ to the nearest point --- real \emph{or complex} ---
where the function misbehaves. The interval on the real line is the diameter of a
hidden disk in the complex plane. So when a series stops converging at a point where
the function looks fine, look for the nearest singularity off the real axis.
\end{connectionbox}

\subsection{A counterexample: smooth but not analytic}

Recall Cauchy's flat function from Chapter~30. It separates
``has a Taylor series'' and ``equals its Taylor series'':
\[
  f(x) =
  \begin{cases}
    e^{-1/x^2}, & x \ne 0, \\
    0, & x = 0.
  \end{cases}
\]
This function is infinitely differentiable everywhere, including at $0$. But it is so
flat at the origin that (as you will verify for $f'(0)$ in the exercises)
\[
  f(0) = f'(0) = f''(0) = \cdots = f^{(n)}(0) = \cdots = 0
\]
--- \emph{every} derivative at $0$ is zero. The exponential $e^{-1/x^2}$ decays to
zero faster than any power of $x$ can detect, so the difference quotients defining
each derivative all collapse to $0$.

Now build the Maclaurin series. Every coefficient $\frac{f^{(k)}(0)}{k!}$ is $0$, so
the Taylor series is
\[
  0 + 0\cdot x + 0\cdot x^2 + \cdots = 0.
\]
The series converges, but to the zero function, and $f(x)$ is \emph{not}
zero for $x \ne 0$. The Taylor series converges everywhere and yet equals $f$ only at
the single point $x=0$. Here $R_n(x) = f(x) - 0 = f(x)$ does \emph{not} go to zero;
it stays equal to the whole function for every $n$ (Figure~\ref{ch31:fig-flat}).

\begin{figure}[H]
\centering
\begin{tikzpicture}
\begin{axis}[
  width=11cm, height=6cm,
  axis lines=middle,
  xlabel={$x$}, ylabel={$y$},
  xmin=-2.6, xmax=2.6, ymin=-0.15, ymax=1.05,
  samples=200, domain=-2.5:2.5,
  clip=true,
]
\addplot[curveblue, very thick] {exp(-1/(x^2+0.0000001))};
\addplot[tangentred, thick, dashed, domain=-2.5:2.5] {0};
\addplot[only marks, mark=*, mark size=1.5pt, black] coordinates {(0,0)};
\node[curveblue] at (axis cs:1.9,0.85) {$e^{-1/x^2}$};
\node[tangentred, anchor=south west] at (axis cs:-2.5,0.03) {Taylor series $= 0$};
\end{axis}
\end{tikzpicture}
\caption{The function $e^{-1/x^2}$ (extended by $0$ at the origin) is infinitely
flat at $x=0$: every derivative there is zero. Its Taylor series is the zero function
(red), which matches the function only at the single point $x=0$. Smoothness does not
guarantee that a function equals its Taylor series.}
\label{ch31:fig-flat}
\end{figure}

\begin{focusbox}
\textbf{Analytic functions.} A function is \emph{analytic} at $a$ if its Taylor
series centered at $a$ converges to the function on some interval around $a$ --- that
is, if $R_n(x) \to 0$ there. Most familiar functions ($e^x$, $\sin$, $\cos$,
polynomials, $\ln$, $1/(1-x)$ away from their singularities) are analytic. The
example $e^{-1/x^2}$ shows that \emph{infinitely differentiable is strictly weaker
than analytic}: being smooth is necessary but not sufficient. Equality with the
Taylor series is an additional property, and $R_n \to 0$ is the exact condition that
certifies it.
\end{focusbox}

\begin{aestheticbox}
Analytic functions are far more rigid than functions that are merely smooth (also written
$C^\infty$, meaning ``derivatives of every order exist and are continuous''). The function
that equals $0$ for $x \le 0$ and $e^{-1/x^2}$ for $x > 0$ is smooth: joining the two pieces at
$x=0$ breaks no derivative. So a smooth function can be zero on one stretch and nonzero on the
next. Analytic functions cannot be patched together this way: knowing all derivatives at a
single point fixes the Taylor series, and therefore the function itself, across the entire
interval of convergence.
\end{aestheticbox}

\section{Physics lives in the first two terms}

Many classic models in physics come down to a single move: keep the first nonvanishing Taylor
term (the first term whose coefficient is not zero, after any constant), drop the higher-order
terms, and solve the resulting equation. The remainder bound says how much this costs.

\subsection{The small-angle approximation}

Models of a swinging pendulum, of light bending, and of a vibrating string all
simplify a great deal when you replace $\sin\theta$ by its first Taylor term. From the series
\[
  \sin\theta = \theta - \frac{\theta^3}{6} + \frac{\theta^5}{120} - \cdots,
\]
for small $\theta$ we keep only the leading term:
\[
  \sin\theta \approx \theta \qquad (\theta \text{ in radians}).
\]
The remainder bound tells you exactly when this is safe. The first dropped term is
$\theta^3/6$, and (using $\abs{f^{(3)}} \le 1$) the Lagrange remainder for the
degree-$1$ approximation is at most $\frac{\abs{\theta}^3}{6}$. At
$\theta = 0.1$ rad (about $5.7^\circ$), the error is below
$\frac{(0.1)^3}{6} \approx 1.7\times 10^{-4}$ --- a relative error well under two tenths
of a percent.

This approximation is behind a famous fact about pendulums. For a pendulum of length
$L$ swinging at angle $\theta(t)$ from the vertical, Newton's laws give
\[
  \theta''(t) = -\frac{g}{L}\sin\theta(t),
\]
where $\theta''$ is the second derivative with respect to time and
$g \approx 9.8\,\text{m/s}^2$ is the acceleration due to gravity. This equation is hard
to solve. Replacing $\sin\theta$ by $\theta$ turns it into $\theta'' = -\frac{g}{L}\theta$,
the equation of a \emph{simple harmonic oscillator} (the same equation describes a mass
bouncing on an ideal spring). Its solutions are sines and cosines with period (the time
for one full swing back and forth)
\[
  T = 2\pi\sqrt{L/g},
\]
which does not depend on the amplitude (how far the pendulum swings). This independence
comes from truncating the Taylor series after the first term, so it holds only for small
swings. Figure~\ref{ch31:fig-small-angle} shows how close $\theta$ and $\sin\theta$ are
near $0$.

\begin{figure}[H]
\centering
\begin{tikzpicture}
\begin{axis}[
  width=11cm, height=6.5cm,
  axis lines=middle,
  xlabel={$\theta$}, ylabel={},
  xmin=-1.6, xmax=1.6, ymin=-1.4, ymax=1.4,
  samples=160, domain=-1.55:1.55,
  legend pos=south east, legend cell align=left,
  clip=true,
]
\addplot[curveblue, very thick] {sin(deg(x))};
\addlegendentry{$\sin\theta$}
\addplot[tangentred, thick, dashed] {x};
\addlegendentry{$\theta$ (small-angle)}
\end{axis}
\end{tikzpicture}
\caption{The small-angle approximation $\sin\theta \approx \theta$ ($\theta$ in radians). Near the origin
the line $y=\theta$ and the curve $y=\sin\theta$ are nearly indistinguishable; the gap
grows like $\theta^3/6$, so it stays tiny while $\theta$ is small.}
\label{ch31:fig-small-angle}
\end{figure}

\begin{applicationbox}
The same ``first nonvanishing term'' logic runs through physics:
\begin{itemize}
\item \textbf{Relativistic energy.} In Einstein's relativity, the exact kinetic energy
(energy of motion) of a mass $m$ moving at speed $v$ is
\[
  E = mc^2\left(\frac{1}{\sqrt{1-v^2/c^2}} - 1\right),
\]
where $c$ is the speed of light. Expanding $1/\sqrt{1-u}$ with $u = v^2/c^2$ for small
$v/c$ gives
\[
  E \approx \frac12 mv^2 + \frac{3}{8}\,\frac{mv^4}{c^2} + \cdots.
\]
The first term is Newton's $\tfrac12 mv^2$ --- so classical kinetic energy is the
leading Taylor term of the relativistic one. The next term is the first relativistic
correction.
\item \textbf{Potential wells.} Let $V(x)$ be the potential energy of a particle at
position $x$; the force on it is $-V'(x)$. An \emph{equilibrium} is a position $x_0$ where
the force is zero, so $V'(x_0)=0$; it is \emph{stable} if a small push makes the particle
oscillate around $x_0$ instead of running away, as at a minimum of $V$ (the bottom of a
``well''). Near a stable equilibrium the Taylor expansion of a smooth $V$ begins
$V(x) \approx V(x_0) + \tfrac12 V''(x_0)(x-x_0)^2$. The quadratic term is an osculating parabola:
\emph{every} smooth potential well behaves locally like a harmonic oscillator
(Figure~\ref{ch31:fig-potential-well}).
\end{itemize}
In each case the Taylor remainder tells you when the approximation may be trusted
and when the next term must be restored.
\end{applicationbox}

\begin{figure}[H]
\centering
\begin{tikzpicture}
\begin{axis}[
  width=11.5cm, height=6.0cm,
  axis lines=middle,
  xlabel={$x$}, ylabel={$V(x)$},
  xmin=-1.8, xmax=2.2, ymin=-1.5, ymax=1.8,
  xtick={-1,1},
  ytick={1},
  samples=120, domain=-1.5:1.9,
  legend style={at={(0.5,-0.08)}, anchor=north, legend columns=2, font=\small, draw=gray!40, rounded corners=2pt},
  clip=true,
]
\addplot[curveblue, very thick, domain=-1.2:1.8] {0.8*(x-0.3)^2 - 0.4*(x-0.3)^3 - 0.7};
\addlegendentry{Potential $V(x)$}
\addplot[tangentred, thick, dashed, domain=-0.8:1.4] {0.8*(x-0.3)^2 - 0.7};
\addlegendentry{Quadratic spring $P_2(x)$}
\addplot[only marks, mark=*, mark size=2pt, black] coordinates {(0.3,-0.7)};
\node[anchor=west, font=\footnotesize] at (axis cs:0.45,-0.82) {equilibrium $x_0$};
\draw[<->, >=stealth, thick, quotientorange] (axis cs:-0.2,-1.1) -- (axis cs:0.8,-1.1)
  node[right=1pt, font=\footnotesize] {small oscillations};
\end{axis}
\end{tikzpicture}
\caption{Near any stable equilibrium $x_0$, the potential
$V(x)$ has $V'(x_0)=0$. The leading non-constant Taylor term is quadratic:
$V(x) \approx V(x_0) + \tfrac12 V''(x_0)(x-x_0)^2$. To a particle confined to small oscillations
around the trough, the true potential is indistinguishable from a parabola --- an effective
Hooke's-law spring with constant $k = V''(x_0)$.}
\label{ch31:fig-potential-well}
\end{figure}

\begin{walkbox}
Walk and turn this over: why is it the \emph{second} derivative that governs the
behavior of a system near equilibrium, never the first? (Hint: equilibrium means the
first derivative of the potential is zero, so the first surviving Taylor term is the
quadratic one --- and a quadratic potential is exactly a spring. Stability is
$V''>0$; instability is $V''<0$.) What would a system look like if $V''$ also
vanished and the leading term were quartic?
\end{walkbox}

\section{Worked examples, start to finish}

\begin{example}[Estimating $\sqrt{e}$ with a guarantee]
\label{ch31:ex-sqrt-e}
Approximate $\sqrt{e} = e^{0.5}$ using the degree-$3$ Maclaurin polynomial of $e^x$
and bound the error.

The polynomial is $P_3(x) = 1 + x + \tfrac{x^2}{2} + \tfrac{x^3}{6}$. At $x=0.5$,
\[
  P_3(0.5) = 1 + 0.5 + \tfrac{0.25}{2} + \tfrac{0.125}{6}
           = 1 + 0.5 + 0.125 + 0.0208\overline{3} = 1.6458\overline{3}.
\]
For the error, $f^{(4)}(x) = e^x$ and on $[0,0.5]$ we have
$e^t \le e^{0.5} < 1.65$, so take $M = 1.65$. Then
\[
  \abs{R_3(0.5)} \le \frac{1.65}{4!}(0.5)^4
  = \frac{1.65}{24}\cdot 0.0625 \approx 0.0043.
\]
So $\sqrt{e} \in [1.6458 - 0.0043,\; 1.6458 + 0.0043] = [1.6415,\, 1.6501]$. The true
value is $\sqrt{e} = 1.64872\ldots$, inside the interval, with actual error
$\approx 0.0029$ --- under the promised $0.0043$.
\end{example}

\begin{example}[A Taylor expansion not centered at $0$]
\label{ch31:ex-ln}
Approximate $\ln(1.2)$ by expanding $\ln x$ about $a = 1$.

The derivatives of $f(x) = \ln x$ are $f'(x) = x^{-1}$, $f''(x) = -x^{-2}$,
$f'''(x) = 2x^{-3}$, so at $a=1$: $f(1)=0$, $f'(1)=1$, $f''(1)=-1$, $f'''(1)=2$. The
degree-$2$ polynomial is
\[
  P_2(x) = 0 + 1\cdot(x-1) + \frac{-1}{2}(x-1)^2 = (x-1) - \tfrac12(x-1)^2.
\]
At $x = 1.2$, with $(x-1) = 0.2$:
\[
  P_2(1.2) = 0.2 - \tfrac12(0.04) = 0.2 - 0.02 = 0.18.
\]
For the error, $f'''(x) = 2x^{-3}$ and on $[1, 1.2]$ this is largest at $x=1$, where
it equals $2$, so $M = 2$. Then
\[
  \abs{R_2(1.2)} \le \frac{2}{3!}(0.2)^3 = \frac{2}{6}(0.008) \approx 0.00267.
\]
The true value $\ln(1.2) = 0.18232\ldots$ differs from $0.18$ by $0.00232$, inside
the bound. We expanded about $1$, not $0$, because $\ln$ is undefined at $0$ ---
you center where the function is friendly and where you know the derivatives.
\end{example}

\begin{example}[Reading off a limit from the first surviving term]
\label{ch31:ex-limit}
Evaluate the following limit using Taylor expansion:
\[
  \lim_{x\to 0}\frac{x - \sin x}{x^3}.
\]

Expand the numerator: $\sin x = x - \tfrac{x^3}{6} + \tfrac{x^5}{120} - \cdots$, so
\[
  x - \sin x = \frac{x^3}{6} - \frac{x^5}{120} + \cdots.
\]
Divide by $x^3$:
\[
  \frac{x - \sin x}{x^3} = \frac{1}{6} - \frac{x^2}{120} + \cdots
  \;\xrightarrow[x\to 0]{}\; \frac{1}{6}.
\]
The limit is $\tfrac16$. This is the same answer L'H\^opital's rule gives after three
differentiations, but Taylor expansion shows why: the
numerator's first nonzero term is $x^3/6$, and dividing by $x^3$ exposes its
coefficient. Taylor series turn many delicate limits into bookkeeping of leading
terms.
\end{example}

\begin{connectionbox}
Example~\ref{ch31:ex-limit} illustrates a general method. Whenever you face a $0/0$ limit,
expanding numerator and denominator in Taylor series and comparing their lowest surviving
powers usually resolves it, and shows why L'H\^opital's rule works.
L'H\^opital differentiates to remove the vanishing terms one order at a time; Taylor expansion
shows all orders at once. Near a point where it is zero, a function behaves like its
lowest-degree nonzero Taylor term, provided it has one (the flat function $e^{-1/x^2}$ does not).
\end{connectionbox}

\section{What to carry away}

Taylor's theorem turns polynomial approximation from a qualitative picture into a quantitative
tool with explicit error bounds.

First, \textbf{Taylor polynomials match derivatives}. The degree-$n$ Taylor
polynomial
\[
  P_n(x) = \sum_{k=0}^n \frac{f^{(k)}(a)}{k!}(x-a)^k
\]
is the unique polynomial of
degree $n$ matching the value and the first $n$ derivatives of $f$ at the anchor point $a$.
The factorials in the denominators cancel the factorials produced by repeated differentiation.

Second, \textbf{Taylor's theorem makes the gap exact}. By writing $f(x) = P_n(x) + R_n(x)$,
we convert approximation into an exact identity. The Lagrange form of the remainder,
\[
  R_n(x) = \frac{f^{(n+1)}(c)}{(n+1)!}\,(x-a)^{n+1},
\]
shows that the error has the form of the next Taylor term, evaluated at an intermediate point
$c$ between $a$ and $x$. This generalizes the Mean Value Theorem: where the MVT balances
average slope with instantaneous slope through one point, Taylor's theorem balances a polynomial
approximation with the curve through the $(n+1)$-st derivative.

Third, \textbf{practical bounds require only a ceiling, not the mystery point $c$}.
The error bound $\abs{R_n(x)} \le \frac{M}{(n+1)!}\abs{x-a}^{n+1}$ demands only an upper bound
$M$ on the next derivative over the interval. When the derivatives stay uniformly bounded,
as for $e^x$, $\sin x$, and $\cos x$ on any fixed interval, the factorial in the denominator
outgrows the power $\abs{x-a}^{n+1}$ and additional terms drive the error to zero.
Bounds of this kind are how calculators guarantee their digits.

Fourth, \textbf{smoothness does not guarantee that the series equals the function}. A Taylor series
equals the function it was built from if and only if the remainder vanishes: $\lim_{n\to\infty} R_n(x) = 0$.
For rational functions like $\frac{1}{1+x^2}$, convergence halts at a finite radius because of
singularities hiding off the real axis in the complex plane. And for Cauchy's flat function
$e^{-1/x^2}$, every derivative at the origin is zero, so the series is identically zero and
tells us nothing about the function away from the origin. Only \emph{analytic} functions are
locally determined by their Taylor series.

\subsection{Looking forward: multivariable calculus}

Up to this point, calculus has operated in one dimension: functions $y=f(x)$ took a single input
and returned a single output. In Chapter~32, we move
to functions of several variables, where inputs are points in the plane or in space, graphs become
surfaces, and derivatives measure rates of change along multiple directions.

\section{Exercises}

\begin{exercise}
Find the degree-$3$ Maclaurin polynomial of $f(x) = e^{x}\cos x$ by computing
$f(0)$, $f'(0)$, $f''(0)$, and $f'''(0)$ directly. Confirm that the result matches
$1 + x - \tfrac{x^3}{3}$.
\end{exercise}

\begin{exercise}
Use the degree-$2$ Taylor polynomial of $\sqrt{x}$ centered at $a = 4$ to estimate
$\sqrt{4.2}$. Then use the Lagrange remainder to bound the error, and compare with a
calculator value.
\end{exercise}

\begin{exercise}
The small-angle approximation $\sin\theta \approx \theta$ is used for
$\theta = 10^\circ$. Convert to radians, and use the remainder bound
$\abs{R_1(\theta)} \le \abs{\theta}^3/6$ to find the maximum possible error. Is the
approximation good to three decimal places?
\end{exercise}

\begin{exercise}
How many terms of the Maclaurin series for $\cos x$ are needed to compute
$\cos(1)$ with error less than $10^{-8}$? Justify with the Lagrange bound, then state
the polynomial you would use.
\end{exercise}

\begin{exercise}
Use Taylor expansions to evaluate
$\displaystyle \lim_{x\to 0}\frac{e^{x} - 1 - x}{x^2}$. Identify the first surviving
term of the numerator and read off the answer.
\end{exercise}

\begin{exercise}[Conceptual]
A classmate says, ``If a function is infinitely differentiable, its Taylor series must
converge to it.'' Give the counterexample from this chapter and explain in one or two
sentences exactly which step of that reasoning fails.
\end{exercise}

\begin{exercise}[Conceptual]
The Maclaurin series of $\frac{1}{1-x}$ is $1 + x + x^2 + \cdots$, with radius of
convergence $1$. Explain why $1$ is the radius, both from the real-line behavior of
the function and from the location of its singularity. Then predict the radius of
convergence of the Maclaurin series of $\frac{1}{4+x^2}$ without computing the series,
and justify your prediction.
\end{exercise}

\begin{exercise}[Proof]
Prove the integral form of the remainder: if $f^{(n+1)}$ is continuous on the interval
between $a$ and $x$, then
\[
  R_n(x) = \frac{1}{n!}\int_a^{x} f^{(n+1)}(t)\,(x-t)^n\,\dd t.
\]
(Hint: start from $f(x) - f(a) = \int_a^x f'(t)\,\dd t$ and integrate by parts
repeatedly, choosing the antiderivative of $1$ to be $-(x-t)$ so the powers build up.
The Lagrange form follows from this integral form by the
weighted Mean Value Theorem for integrals: if $w$ is continuous and keeps one sign on an
interval and $F$ is continuous there, then $\int F\,w = F(c)\int w$ for some $c$ in the
interval. Apply it with $w(t) = (x-t)^n$.)
\end{exercise}

\begin{exercise}[Proof]
Show that for $f(x) = e^{-1/x^2}$ (with $f(0)=0$), the first derivative at $0$ is
$f'(0) = 0$. (Compute $\lim_{h\to 0}\frac{e^{-1/h^2}}{h}$ by substituting
$u = 1/h^2 \to \infty$ and comparing the growth of $e^{u}$ against $\sqrt{u}$.) Explain
in words why the same kind of argument forces \emph{every} derivative at $0$ to vanish.
\end{exercise}

\begin{exercise}
Expand $f(x) = \frac{1}{\sqrt{1+x}}$ to degree $2$ about $a=0$ to derive the
approximation $\frac{1}{\sqrt{1+x}} \approx 1 - \tfrac12 x + \tfrac38 x^2$. Use it to
estimate $\frac{1}{\sqrt{1.1}}$, bound the error, and explain how this is the Taylor
origin of the relativistic correction term mentioned in the chapter.
\end{exercise}
